\chapter{El Motor de Juego ZINA}

El \textbf{Sistema ZINA} es un reglamento narrativo ligero adaptado de los principios de diseño de \emph{Powered by the Apocalypse} (PbtA) y \emph{The Veil}. En este sistema no existen tiradas pasivas de percepción ni estadísticas físicas convencionales. Lo que impulsa la historia y determina el desenlace de los riesgos es la \textbf{emoción} con la que el personaje actúa en la ficción.

\section{Filosofía de juego: La conversación}

Una partida de NeoBilbao fluye como una conversación entre los jugadores y el Director de Juego:
\begin{enumerate}
  \item El Director describe la situación opresiva, el entorno y los peligros inmediatos.
  \item Los jugadores responden describiendo qué hacen, dicen o intentan sus personajes.
  \item Solo cuando la acción del personaje encaja con el disparador de un \textbf{Movimiento} y conlleva un riesgo real, se toman los dados. La ficción dicta la regla, y el resultado de los dados transforma la ficción.
\end{enumerate}

\section{Las emociones del Sistema ZINA}

En lugar de Fuerza o Inteligencia, los personajes poseen valores emocionales que reflejan su estado psicológico y su forma de relacionarse con el mundo:

\begin{itemize}
  \item \lead{Enfadado:} Furia de clase, frustración acumulada, indignación moral y violencia visceral. Se utiliza cuando el personaje ataca de frente, rompe una puerta a patadas o confronta a un capataz.
  \item \lead{Asustado:} Instinto de supervivencia, adrenalina pura, terror y prudencia. Se emplea cuando el personaje huye a la desesperada, se oculta en las sombras o esquiva un peligro letal.
  \item \lead{Culpable:} El peso de los errores pasados, la vergüenza y el auto-reproche. Es la emoción que impulsa a violar códigos de seguridad, hackear terminales corporativas o asumir riesgos suicidas para enmendar una falta.
  \item \lead{Afligido:} Empatía, dolor por la pérdida de seres queridos, tristeza compartida y compasión. Se utiliza para consolar a un herido, tratar traumas médicos o apelar a los sentimientos de otra persona.
  \item \lead{Pacífico:} Calma forzada, desapego analítico, sangre fría y meticulosidad. Se utiliza cuando el personaje opera maquinaria con pulso firme, falsifica un albarán burocrático o analiza un plano técnico sin alterarse.
  \item \lead{Alegre:} Esperanza activa, humor rebelde, calidez humana y orgullo de barrio.
    \begin{itemize}
      \item \textbf{Activo en 2084:} Representa la fuerza moral de la clase trabajadora de La Margen antes del encierro total.
      \item \textbf{Extinguido en 2102:} En los Zócalos Industriales de 2102, esta emoción está tachada con una \textbf{X} en la hoja de personaje; la alegría se ha perdido tras la Guerra del Plasma y el pacto forzado del Velo.
    \end{itemize}
\end{itemize}

\section{La mecánica de resolución (2d6 + Emoción)}

Cuando un movimiento requiere tirar dados, el jugador lanza \textbf{dos dados de seis caras (2d6)} y suma el modificador de la emoción que mejor describe su actitud:

\begin{CajaRegla}{La escala de tres resultados de ZINA}
\begin{itemize}
  \item \textbf{10 o más (Éxito Total):} Logras exactamente lo que pretendías. Tomas la iniciativa en la escena y no sufres consecuencias negativas imprevistas.
  \item \textbf{7 a 9 (Éxito con Coste):} Consigues tu objetivo básico, pero el Director introduce una complicación, un dilema o un precio: sufres daño, aumentas tu nivel de Zama, pierdes un objeto valioso o alertas a patrullas cercanas.
  \item \textbf{6 o menos (Fallo):} Las cosas se tuercen drásticamente. El Director toma el control y realiza un \textbf{Movimiento Duro}: una emboscada, una sobrecarga en el reactor, la rotura de un equipo vital o un aumento crítico de la tensión.
\end{itemize}
\end{CajaRegla}

\section{Zama (Carga y Colapso)}

La \textbf{Zama} representa el agotamiento mental, la acumulación de neurotoxinas y el desgaste psicológico que soporta el personaje.
\begin{itemize}
  \item La escala va de \textbf{0 a 6 puntos}.
  \item Todo personaje comienza con \textbf{0 de Zama}.
  \item El estrés aumenta al inhalar gases ácidos, al fallar tiradas traumáticas o al sufrir ataques neurotóxicos de los Auditores o la División Sorgin.
  \item \lead{Alcanzar 6 de Zama (El Colapso):} El personaje no muere necesariamente, pero queda incapacitado. Puede sufrir un ataque de pánico ciego, caer en coma catatónico, desmayarse por sobrecarga o acatar sumisamente las órdenes corporativas por terror psicológico.
  \item \lead{Aliviar la Zama:} Se reduce compartiendo descansos en refugios seguros, consumiendo fármacos de urgencia o desahogándose fraternalmente en una tasca.
\end{itemize}

\section{La Reserva de Auzolan (Solidaridad)}

El \textbf{Auzolan} es el fondo compartido de apoyo obrero y vecinal:
\begin{itemize}
  \item Al inicio de cada sesión de juego, el grupo recibe \textbf{1 punto de Auzolan por cada jugador en la mesa}.
  \item Cualquier jugador puede gastar puntos de la reserva en cualquier momento para:
    \begin{enumerate}
      \item \textbf{Echar un cable (+1):} Otorgar un bono de +1 a la tirada de un compañero cercano tras describir cómo le ayuda físicamente o le grita una advertencia.
      \item \textbf{Declarar un recurso de barrio:} Declarar que en la escena existe un contacto conocido, una trampilla de escape o una caja de herramientas olvidada.
      \item \textbf{Absorber el golpe:} Reducir en 1 punto el daño o la Zama que fuera a recibir un aliado adyacente.
      \item \textbf{Pausa comunal:} Durante un descanso seguro en una tasca o barracón, gastar 1 punto de Auzolan reduce en 1 punto la Zama de todos los personajes presentes.
    \end{enumerate}
\end{itemize}

\section{Movimientos Básicos Comunes}

Todos los personajes, independientemente de su arquetipo, tienen acceso a los siguientes movimientos fundamentales:

\subsection{1. Atacar / Forcejear (+Enfadado)}
Cuando recurres a la fuerza bruta, embistes a un guardia o golpeas con una herramienta pesada:
\begin{itemize}
  \item \textbf{10+:} Derribas o inutilizas a tu adversario sin sufrir daño a cambio.
  \item \textbf{7--9:} Logras herir o apartar a tu oponente, pero él te alcanza en el forcejeo: sufres daño físico o aumentas tu Zama en 1 punto.
  \item \textbf{Fallo:} Tu ataque falla estrepitosamente o tu arma se atasca; quedas a merced del contraataque.
\end{itemize}

\subsection{2. Infiltrarse / Ocultarse (+Asustado)}
Cuando te mueves entre las sombras, esquivas el haz de un dron Tártalo o te escondes tras un contenedor:
\begin{itemize}
  \item \textbf{10+:} Pasas totalmente desapercibido y alcanzas tu posición de cobertura.
  \item \textbf{7--9:} Llegas a tu destino, pero dejas una pista visible, haces ruido o te ves obligado a abandonar un objeto de tu equipo para no ser visto.
  \item \textbf{Fallo:} Quedas acorralado en un callejón sin salida bajo el foco de seguridad.
\end{itemize}

\subsection{3. Hackear / Puentear Sistemas (+Culpable)}
Cuando conectas cables a una terminal, descifras un código de esclusa o sobrecargas un panel eléctrico:
\begin{itemize}
  \item \textbf{10+:} Obtienes acceso total, extraes los datos o activas el mecanismo limpiamente.
  \item \textbf{7--9:} Consigues el acceso, pero el cortocircuito activa una alarma silenciosa o quema un componente de tu terminal portátil.
  \item \textbf{Fallo:} La terminal se bloquea, emite una descarga eléctrica (1 de Zama) y transmite tu ubicación al puesto de control.
\end{itemize}

\subsection{4. Manipular / Conectar (+Afligido)}
Cuando apelas al dolor compartido, consuelas a una víctima o intentas conmover a un trabajador para que colabore:
\begin{itemize}
  \item \textbf{10+:} La persona empatiza contigo, se abre y te concede la ayuda o información solicitada.
  \item \textbf{7--9:} Colabora, pero exige una garantía a cambio: comida, medicinas o que asumas un riesgo por ella.
  \item \textbf{Fallo:} La persona se cierra en banda, dominada por el miedo a las represalias corporativas, o te delata a la guardia.
\end{itemize}

\subsection{5. Analizar con Sangre Fría (+Pacífico)}
Cuando estudias un plano arquitectónico, examinas una herida bajo presión o intentas mantener la compostura ante un interrogatorio:
\begin{itemize}
  \item \textbf{10+:} El Director te revela con claridad la mayor debilidad de la escena o la mejor ruta de escape. Tienes +1 a tu siguiente acción al actuar según esa información.
  \item \textbf{7--9:} Descubres un detalle crucial, pero pierdes segundos valiosos o te expones a ser sorprendido.
  \item \textbf{Fallo:} Tu análisis es erróneo o te bloqueas mentalmente por la tensión.
\end{itemize}

\subsection{6. Infundir Esperanza (+Alegre -- Solo Era 2084)}
Cuando cuentas un chiste en mitad del peligro, recuerdas a tus compañeros que saldréis de esta o cantas un himno obrero en la tasca:
\begin{itemize}
  \item \textbf{10+:} Tus compañeros recuperan la calma: tú y un aliado adyacente reducís 1 punto de Zama inmediatamente.
  \item \textbf{7--9:} Levantas la moral del grupo, pero tu optimismo llama la atención de quien no debe o te hace bajar la guardia.
  \item \textbf{Fallo:} La realidad aplastante apaga tus palabras; el grupo sufre 1 punto de Zama por desesperación.
\end{itemize}
